\section{Trees (Árvores Binárias)}
\label{sec:modeling-trees}

A categoria \textbf{Trees} é composta por árvores binárias finitas como objetos, onde a estrutura é definida pela hierarquia entre nós e pela lateralidade das ramificações. Esta categoria é essencial para modelar processos de decisão dicotómicos e estruturas de dados hierárquicas onde a distinção entre esquerda e direita é semanticamente relevante.

\subsection{Objetos}
No ambiente MATLAB/Octave, os objetos desta categoria podem ser representados de duas formas complementares, dependendo da operação ou construção categorial a realizar:

\paragraph{1. Representação por Array de Progenitores:}
Esta forma utiliza um vetor \texttt{parent} (onde \texttt{parent(i)} é o índice do pai do nó $i$) e um vetor \texttt{side} (onde $0$ representa o filho esquerdo e $1$ o filho direito). É a representação mais eficiente para validar propriedades de morfismos e calcular limites.

\begin{lstlisting}[language=Matlab, caption={Representação por Array de Progenitores}]
% Arvore com 7 nos: raiz=1, filhos esquerdos/direitos indicados
%       1
%      / \
%     2   3
%    / \ / \
%   4  5 6  7

n      = 7;
parent = [0, 1, 1, 2, 2, 3, 3];  % parent(1)=0 -> raiz
side   = [-1, 0, 1, 0, 1, 0, 1]; % -1=raiz, 0=esq, 1=dir

tree.n      = n;
tree.parent = parent;
tree.side   = side;
\end{lstlisting}

\paragraph{2. Representação Recursiva:}
Tal como explorado no script original, esta representação define a árvore como uma estrutura (\textit{struct}) que contém um valor e referências para as suas subárvores. Esta abordagem é fundamental para a definição de construções como o Produto de árvores, facilitando a navegação descendente.

\begin{lstlisting}[language=Matlab, caption={Representação Recursiva de uma Árvore Binária}]
% Definicao de um no com subarvores
node.value = 10;    % Conteudo informativo do no
node.left  = [];    % Subarvore esquerda (vazia se [])
node.right = [];    % Subarvore direita (vazia se [])

% Exemplo de construcao de uma arvore de altura 1
T.value = 1;
T.left  = struct('value', 2, 'left', [], 'right', []);
T.right = struct('value', 3, 'left', [], 'right', []);
\end{lstlisting}

\subsection{Morfismos}
Os morfismos $f: T_1 \to T_2$ nesta categoria são mapeamentos entre os conjuntos de nós que preservam estritamente a raiz, a relação de parentesco e a lateralidade. Para que $f$ seja um morfismo válido, a imagem do pai de um nó em $T_1$ deve coincidir com o pai da imagem desse nó em $T_2$, mantendo-se a posição relativa (lado).

\begin{lstlisting}[language=Matlab, caption={Algoritmo de Verificação de Morfismo de Árvores}]
function ok = is_tree_morphism(T1, T2, f)
    ok = true;
    
    % 1) Verificacao da Raiz: a imagem da raiz deve ser a raiz
    r1 = find(T1.parent == 0);
    r2 = find(T2.parent == 0);
    if f(r1) ~= r2, ok = false; return; end
    
    % 2) Preservacao de Estrutura e Lateralidade
    for i = 1:T1.n
        if T1.parent(i) ~= 0
            p1 = T1.parent(i);
            % O pai da imagem deve ser a imagem do pai
            if T2.parent(f(i)) ~= f(p1), ok = false; return; end
            % O lado (esquerdo/direito) deve ser preservado
            if T2.side(f(i)) ~= T1.side(i), ok = false; return; end
        end
    end
end
\end{lstlisting}